%=============================================================================
% Verification/Falsification: The Electroweak Derivation kappa = cos 2theta_W
% Critical analysis of the DFD prediction for the constitutive split parameter
% March 2026
%=============================================================================

\documentclass[aps,prd,twocolumn,superscriptaddress,nofootinbib]{revtex4-2}

\usepackage{amsmath,amssymb,amsfonts}
\usepackage{hyperref}
\usepackage{booktabs}
\usepackage{tcolorbox}

\newcommand{\CP}{\mathbb{C}P}
\newcommand{\Tr}{\operatorname{Tr}}
\newcommand{\vE}{\mathbf{E}}
\newcommand{\vB}{\mathbf{B}}

\begin{document}

\title{Critical Verification: $\kappa = \cos 2\theta_W \approx 0.54$\\
The Electroweak Derivation of the Constitutive Split Parameter}

\author{DFD Research Programme}

\date{23 March 2026}

\begin{abstract}
We subject the DFD claim $\kappa = \cos 2\theta_W$ to systematic
verification.  Five independent checks are performed: (1)~whether the
spectral action on $\CP^2 \times S^3$ couples $\psi$ asymmetrically
to $U(1)_Y$ vs.\ $SU(2)_L$; (2)~whether the Frame Stiffness Theorem
ratio $\kappa_{U(1)}/\kappa_{SU(2)} = 1/2$ determines $\kappa$;
(3)~what observational consequences follow from symmetric ($\kappa = 0$)
vs.\ asymmetric loading; (4)~whether the formula
$\cos 2\theta_W = 1 - 2\sin^2\theta_W$ yields $7/13$ with the DFD
prediction $\sin^2\theta_W = 3/13$; (5)~whether the derivation uses
assumptions beyond the DFD framework.  We find the derivation
\textbf{logically valid within DFD} but identify one critical
assumption that requires independent justification and flag a
tension with an alternative $\kappa$ route.
\end{abstract}

\maketitle


%=============================================================================
\section{The Claim Under Review}
\label{sec:claim}
%=============================================================================

The constitutive relations of the DFD loaded vacuum are
\begin{equation}
\epsilon(\psi) = \epsilon_0\, n\, e^{+\kappa\psi}, \qquad
\mu_{\mathrm{mag}}(\psi) = \mu_0\, n\, e^{-\kappa\psi},
\end{equation}
with $n = e^\psi$.  The parameter $\kappa$ controls the
electric/magnetic asymmetry of the vacuum response to the density
field $\psi$.

The claim (Vacuum Loading Deep Physics, Sec.~2.4, Eq.~(14)):
\begin{equation}
\boxed{\kappa = \cos 2\theta_W = 1 - 2\sin^2\theta_W \approx 0.538.}
\label{eq:claim}
\end{equation}

The physical argument: the photon is a mixture
$A_\mu^{\mathrm{EM}} = \sin\theta_W\, W^3_\mu + \cos\theta_W\, B_\mu$.
The $\psi$-field couples asymmetrically to these sectors because $\psi$
couples to the $U(1)_Y$ hypercharge sector (through the spectral action
on $\CP^2 \times S^3$) but not directly to $SU(2)_L$ (which lives on a
different part of the internal geometry).  Therefore $\kappa = \cos 2\theta_W$.


%=============================================================================
\section{Check 1: Does the Spectral Action Couple $\psi$ Asymmetrically?}
\label{sec:check1}
%=============================================================================

\subsection{The internal geometry}

In DFD gauge emergence, the internal space is $\CP^2 \times S^3$.
The gauge sectors correspond to isometries acting on subspaces:
\begin{center}
\begin{tabular}{cccc}
\toprule
Sector & Subspace & Dim & Ricci factor \\
\midrule
$SU(3)$ & $\CP^2$ & 3 & 3 \\
$SU(2)$ & $\CP^1 \subset \CP^2$ & 2 & 2 \\
$U(1)_Y$ & $\CP^0$ (point) & 1 & 1 \\
\bottomrule
\end{tabular}
\end{center}

The Frame Stiffness Theorem (Appendix~F, Theorem~F.7) gives:
\begin{equation}
\kappa_r = n_r \kappa_0, \qquad
\kappa_{SU(3)} : \kappa_{SU(2)} : \kappa_{U(1)} = 3 : 2 : 1.
\label{eq:stiffness}
\end{equation}

\subsection{How $\psi$ couples to gauge sectors}

The DFD scalar $\psi$ couples to gauge fields through the optical metric
$\tilde{g}_{\mu\nu} = e^{2\psi}\eta_{\mu\nu}$.  This coupling is
\emph{universal}---it modifies the effective metric seen by all gauge
bosons.  However, the \emph{stiffness} of each gauge sector against
$\psi$ perturbations differs because the Ricci curvature of the
corresponding projective space differs.

The $U(1)_Y$ sector lives on $\CP^0$ (stiffness $\kappa_0$).  The
$SU(2)_L$ sector lives on $\CP^1$ (stiffness $2\kappa_0$).  A stiffer
sector resists $\psi$-deformation more strongly.

\begin{tcolorbox}[colback=green!5!white, colframe=green!60!black, title=Verdict: Check 1]
\textbf{VERIFIED.}  The spectral action does couple $\psi$
asymmetrically to $U(1)_Y$ vs.\ $SU(2)_L$.  The asymmetry is
quantified by the stiffness ratio $\kappa_{U(1)}/\kappa_{SU(2)} = 1/2$.
The $U(1)_Y$ sector is softer (more responsive to $\psi$) than
$SU(2)_L$, which is the correct direction for $\kappa > 0$
(electric sector more responsive than magnetic).
\end{tcolorbox}


%=============================================================================
\section{Check 2: Does the Stiffness Ratio $1/2$ Determine $\kappa$?}
\label{sec:check2}
%=============================================================================

\subsection{From stiffness to constitutive split}

The photon is a linear combination of $B_\mu$ (from $U(1)_Y$) and
$W^3_\mu$ (from $SU(2)_L$):
\begin{equation}
A_\mu^{\mathrm{EM}} = \cos\theta_W\, B_\mu + \sin\theta_W\, W^3_\mu.
\end{equation}

When $\psi$ perturbs the vacuum, the $B$-component responds with
effective stiffness $\kappa_{U(1)} = \kappa_0$ and the $W^3$-component
with stiffness $\kappa_{SU(2)} = 2\kappa_0$.  The photon's net
$\psi$-response is the weighted average:
\begin{equation}
\delta A_\mu^{\mathrm{EM}} \propto
\frac{\cos^2\theta_W}{\kappa_{U(1)}} B_\mu
+ \frac{\sin^2\theta_W}{\kappa_{SU(2)}} W^3_\mu.
\end{equation}

The constitutive parameter $\kappa$ measures the electric/magnetic
asymmetry.  The $B$-boson (hypercharge) contributes to the
\emph{electric} response with weight $\cos^2\theta_W$, while $W^3$
contributes with weight $\sin^2\theta_W$.

\subsection{The argument for $\cos 2\theta_W$}

The Vacuum Loading paper presents the derivation as:
\begin{equation}
\kappa = \frac{\cos^2\theta_W - \sin^2\theta_W}
{\cos^2\theta_W + \sin^2\theta_W} = \cos 2\theta_W.
\end{equation}

This is a \emph{mixing-fraction} argument: the $B$-component
(fraction $\cos^2\theta_W$ of the photon) couples more strongly to
$\psi$ than the $W^3$-component (fraction $\sin^2\theta_W$).  The
asymmetry is the difference divided by the sum.

\subsection{Critical assessment: is this the right weighting?}

There is a subtlety here.  The formula $\kappa = \cos 2\theta_W$
treats the two gauge boson components as contributing with equal
intrinsic coupling to $\psi$, differing only in their mixing fraction.
But the Frame Stiffness Theorem says the $U(1)_Y$ sector has
stiffness $\kappa_0$ while $SU(2)_L$ has stiffness $2\kappa_0$.

If we incorporate the stiffness asymmetry \emph{as well as} the
mixing, the $\psi$-coupling of each component becomes:
\begin{align}
\text{$B$-component:}& \quad \cos^2\theta_W \times \frac{1}{\kappa_{U(1)}}
= \frac{\cos^2\theta_W}{\kappa_0}, \\
\text{$W^3$-component:}& \quad \sin^2\theta_W \times \frac{1}{\kappa_{SU(2)}}
= \frac{\sin^2\theta_W}{2\kappa_0}.
\end{align}

The asymmetry parameter would then be:
\begin{equation}
\kappa_{\mathrm{stiff}} =
\frac{\cos^2\theta_W/\kappa_0 - \sin^2\theta_W/(2\kappa_0)}
{\cos^2\theta_W/\kappa_0 + \sin^2\theta_W/(2\kappa_0)}
= \frac{2\cos^2\theta_W - \sin^2\theta_W}
{2\cos^2\theta_W + \sin^2\theta_W}.
\label{eq:kappa-stiff}
\end{equation}

With $\sin^2\theta_W = 3/13$:
\begin{equation}
\kappa_{\mathrm{stiff}} = \frac{2(10/13) - 3/13}{2(10/13) + 3/13}
= \frac{20/13 - 3/13}{20/13 + 3/13}
= \frac{17}{23} \approx 0.739.
\end{equation}

This differs from $\cos 2\theta_W = 7/13 \approx 0.538$.

\begin{tcolorbox}[colback=yellow!5!white, colframe=yellow!60!black, title=Verdict: Check 2]
\textbf{TENSION IDENTIFIED.}  The stiffness ratio $1/2$ and the
mixing angle are both relevant, but they enter differently depending
on whether one treats the stiffness as modifying the intrinsic
$\psi$-coupling (giving $\kappa \approx 0.74$) or treats the
mixing as the sole source of asymmetry (giving
$\kappa = \cos 2\theta_W \approx 0.54$).

The Vacuum Loading paper uses the pure mixing formula.  This is
correct \emph{if and only if} the $\psi$-coupling to the photon
is mediated through the electromagnetic field strength $F_{\mu\nu}$
(which is already the physical mixture), rather than through the
separate $B$ and $W^3$ field strengths.  Since DFD couples $\psi$
to the \emph{optical metric} $\tilde{g}_{\mu\nu} = e^{2\psi}\eta_{\mu\nu}$,
which modifies the photon propagation as a whole, the pure mixing
formula is the more natural interpretation.

\textbf{Resolution:} The stiffness ratio determines the mixing
angle $\sin^2\theta_W = 3/13$, and then the mixing angle determines
$\kappa$.  These are sequential, not simultaneous: the stiffness
enters at the level of coupling constants, and $\kappa$ is
determined at the level of the photon's composition.  The pure
mixing formula $\kappa = \cos 2\theta_W$ is therefore
\textbf{self-consistent}, provided the $\psi$-coupling is to the
physical photon field.
\end{tcolorbox}


%=============================================================================
\section{Check 3: Symmetric ($\kappa = 0$) vs.\ Asymmetric Loading}
\label{sec:check3}
%=============================================================================

\subsection{If $\kappa = 0$ (symmetric loading)}

With $\kappa = 0$, the constitutive relations become:
\begin{equation}
\epsilon(\psi) = \epsilon_0\, e^\psi, \qquad
\mu_{\mathrm{mag}}(\psi) = \mu_0\, e^\psi.
\end{equation}
Both electric and magnetic sectors shift identically.  Consequences:
\begin{enumerate}
\item \textbf{No TE/TM splitting:} The TE and TM modes of a cavity
experience identical frequency shifts under $\psi$.  The predicted
splitting $\delta\nu_{\mathrm{TE}} - \delta\nu_{\mathrm{TM}}$
vanishes identically.

\item \textbf{Impedance invariance:} The vacuum impedance
$Z_0 = \sqrt{\mu/\epsilon} = \sqrt{\mu_0/\epsilon_0}$ is independent
of $\psi$.  No gravitational birefringence or impedance variation.

\item \textbf{LPI test insensitivity:} The Local Position Invariance
violation parameter $\xi$ becomes insensitive to the $\psi$-field
in the electromagnetic sector.  Clock comparison experiments testing
$\xi$ at the $10^{-7}$ level would see no signal.

\item \textbf{Electroweak inconsistency:} If $\psi$ couples
symmetrically to $U(1)_Y$ and $SU(2)_L$, the stiffness ratio would
be $1:1$, contradicting the Frame Stiffness Theorem $\kappa_r = n_r \kappa_0$.
\end{enumerate}

\subsection{If $\kappa = \cos 2\theta_W \approx 0.54$}

\begin{enumerate}
\item \textbf{TE/TM splitting:}
$(\delta\nu_{\mathrm{TE}} - \delta\nu_{\mathrm{TM}})/\nu
\approx 2\kappa\psi \approx 4.5 \times 10^{-6}$
at Earth's surface.  This is within reach of SRF cavity technology.

\item \textbf{Modified Coulomb field:} Near the Sun,
$\delta E/E = -(1+\kappa)\psi_\odot \approx -6.5 \times 10^{-6}$.

\item \textbf{Gravitational birefringence:} The vacuum impedance
acquires a $\psi$-dependent correction
$Z(\psi) = Z_0\, e^{-\kappa\psi}$.

\item \textbf{LPI violation:} Measurable at the $10^{-7}$ level
via clock comparisons.
\end{enumerate}

\begin{tcolorbox}[colback=green!5!white, colframe=green!60!black, title=Verdict: Check 3]
\textbf{VERIFIED.}  $\kappa = 0$ is inconsistent with the DFD
framework (contradicts Frame Stiffness Theorem) and produces no
testable EM signatures.  $\kappa \neq 0$ is required by internal
consistency.  The specific value $\kappa = \cos 2\theta_W$ produces
clear, falsifiable predictions.  The key experimental discriminant
is the TE/TM cavity splitting.
\end{tcolorbox}


%=============================================================================
\section{Check 4: Is the Formula Correct? DFD Value $\sin^2\theta_W = 3/13$}
\label{sec:check4}
%=============================================================================

\subsection{Algebraic verification}

The double-angle formula:
\begin{equation}
\cos 2\theta_W = 1 - 2\sin^2\theta_W.
\end{equation}
This is a standard trigonometric identity.  \textbf{Correct.}

\subsection{With the DFD prediction $\sin^2\theta_W = 3/13$}

The DFD derivation (PRL\_SM\_from\_Topology, Appendix~F) gives:
\begin{equation}
\sin^2\theta_W = \frac{g'^2}{g^2 + g'^2}
= \frac{3}{13} = 0.23077\ldots
\end{equation}
derived from:
\begin{itemize}
\item Frame stiffness ratio $\kappa_{SU(2)}/\kappa_{U(1)} = 2$
\item Trace normalization: $\Tr(T^a T^b) = \tfrac{1}{2}\delta^{ab}$
for $SU(2)$; $\Tr(Y^2) = 10/3$ per generation for $U(1)_Y$
\item Combined: $g'^2/g^2 = \kappa_2 \cdot (1/2) / [\kappa_1 \cdot (10/3)] = 3/10$
\item Therefore $\sin^2\theta_W = (3/10)/(1 + 3/10) = 3/13$
\end{itemize}

Substituting into $\kappa$:
\begin{equation}
\kappa = 1 - 2 \times \frac{3}{13} = 1 - \frac{6}{13}
= \frac{7}{13} = 0.53846\ldots
\label{eq:kappa-exact}
\end{equation}

\subsection{Comparison with experimental value}

Using the PDG 2024 value $\sin^2\theta_W^{\overline{\mathrm{MS}}}(M_Z) = 0.23122 \pm 0.00004$:
\begin{equation}
\kappa_{\mathrm{exp}} = 1 - 2(0.23122) = 0.53756 \pm 0.00008.
\end{equation}

The DFD prediction $\kappa = 7/13 = 0.53846$ differs from the
experimental central value by:
\begin{equation}
\frac{|\kappa_{\mathrm{DFD}} - \kappa_{\mathrm{exp}}|}
{\kappa_{\mathrm{exp}}} = \frac{0.00090}{0.53756} = 0.17\%.
\end{equation}

This discrepancy inherits entirely from the offset between
$\sin^2\theta_W = 3/13 = 0.23077$ and the measured $0.23122$.
The DFD value is a tree-level prediction; the 0.2\% offset
is consistent with expected radiative corrections.

\begin{tcolorbox}[colback=green!5!white, colframe=green!60!black, title=Verdict: Check 4]
\textbf{VERIFIED.}  The formula $\cos 2\theta_W = 1 - 2\sin^2\theta_W$
is an exact trigonometric identity.  With $\sin^2\theta_W = 3/13$
(DFD prediction), $\kappa = 7/13 \approx 0.538$.  Using the PDG value,
$\kappa \approx 0.538$.  The two agree to 0.17\%.  The exact DFD
prediction is the rational number $\kappa = 7/13$.
\end{tcolorbox}


%=============================================================================
\section{Check 5: Assumptions Beyond the DFD Framework}
\label{sec:check5}
%=============================================================================

The derivation chain for $\kappa = \cos 2\theta_W$ uses:

\begin{enumerate}
\item \textbf{Internal geometry $\CP^2 \times S^3$.}
Core DFD assumption.  Established by minimality of the $(3,2,1)$ partition
(Appendix~F, Proposition~F.1).

\item \textbf{Frame Stiffness Theorem: $\kappa_r = n_r \kappa_0$.}
Derived from Ricci curvature of projective subspaces
(Appendix~F, Theorem~F.7).  Internal to DFD.

\item \textbf{Gauge coupling ratios from stiffness + trace normalization.}
Combines the stiffness ratio with the standard $\Tr(T^a T^b)$ convention
and the SM hypercharge spectrum $\Tr(Y^2) = 10/3$.  The trace normalization
uses SM field content, which in DFD is derived from the Dirac index on
$\CP^2 \times S^3$.  Internal to DFD.

\item \textbf{$\sin^2\theta_W = 3/13$.}
Follows from items 2--3.  Internal to DFD.

\item \textbf{Photon composition: $A_\mu = \cos\theta_W B_\mu + \sin\theta_W W^3_\mu$.}
Standard electroweak theory.  Not specific to DFD but consistent with it.

\item \textbf{$\kappa$ determined by photon mixing fractions.}
\textbf{This is the critical step.}  The argument is that the constitutive
split $\kappa$ equals the fractional asymmetry
$(\cos^2\theta_W - \sin^2\theta_W)/(\cos^2\theta_W + \sin^2\theta_W)$.
This requires:
\begin{itemize}
\item[(a)] The $\psi$-coupling is to the physical photon as a unit,
  not to the separate $B$ and $W^3$ components individually.
\item[(b)] The $B$-component contributes positively to $\kappa$
  (enhances $\epsilon$ relative to $\mu_{\mathrm{mag}}$) while the
  $W^3$-component contributes negatively (or zero).
\item[(c)] The two contributions enter with equal intrinsic strength,
  differing only in their mixing weights.
\end{itemize}
\end{enumerate}

\subsection{Assessment of assumption 6}

Assumption 6(a) is natural in DFD: the $\psi$ field modifies the
optical metric $\tilde{g}_{\mu\nu} = e^{2\psi}\eta_{\mu\nu}$, which
the photon (not the $B$ or $W^3$ individually) propagates through.

Assumption 6(b) follows from the physical interpretation: the $B$-boson
mediates hypercharge, which is the ``electric'' part of the coupling,
while $W^3$ mediates isospin, which is the ``magnetic'' (non-Abelian)
part.  In DFD, the $U(1)_Y$ sector on $\CP^0$ is softer than $SU(2)_L$
on $\CP^1$, so the electric response dominates, giving $\kappa > 0$.

Assumption 6(c) is the key non-trivial assumption.  As shown in
Check~2 (Sec.~\ref{sec:check2}), incorporating the stiffness ratio
\emph{within} $\kappa$ (rather than only in $\sin^2\theta_W$) would
give $\kappa \approx 0.74$ instead of $0.54$.  The resolution is
that the stiffness ratio has already been ``used up'' in determining
$\sin^2\theta_W = 3/13$---it should not be double-counted.

\begin{tcolorbox}[colback=yellow!5!white, colframe=yellow!60!black, title=Verdict: Check 5]
\textbf{MOSTLY VERIFIED.}  Assumptions 1--5 are internal to DFD.
Assumption~6 is physically motivated and self-consistent, but rests
on the identification of $\kappa$ with the photon mixing asymmetry.
This is \textbf{not derived from first principles}---it is a
physically natural \emph{ansatz}.  The no-double-counting argument
(stiffness determines $\theta_W$, then $\theta_W$ determines $\kappa$)
is logically consistent but would benefit from a rigorous derivation
starting from the gauge--$\psi$ Lagrangian on $\CP^2 \times S^3$.
\end{tcolorbox}


%=============================================================================
\section{Tension with the Hypercharge Route}
\label{sec:tension}
%=============================================================================

The Vacuum Loading paper also presents an alternative $\kappa$
derivation based on the electric/magnetic hypercharge decomposition,
initially giving $\kappa = 1$ (all electric, no magnetic monopoles)
before corrections.  It further cites the Euler-Heisenberg route
giving $\kappa = 3/11 \approx 0.27$.

Three routes give three different values:

\begin{center}
\begin{tabular}{lcc}
\toprule
\textbf{Route} & $\kappa$ & \textbf{Status} \\
\midrule
Hypercharge asymmetry (naive) & 1 & Overcounting \\
Euler-Heisenberg & $3/11 \approx 0.27$ & Wrong regime \\
Electroweak mixing & $7/13 \approx 0.54$ & Preferred \\
\bottomrule
\end{tabular}
\end{center}

The paper correctly identifies the electroweak route as preferred,
arguing that $\psi$ couples to the full electroweak vacuum.
However, the existence of multiple routes with different answers
indicates that a first-principles derivation from the
gauge--$\psi$ action would be valuable for disambiguation.


%=============================================================================
\section{Overall Verdict}
\label{sec:verdict}
%=============================================================================

\begin{tcolorbox}[colback=blue!5!white, colframe=blue!60!black,
  title=Summary Assessment]

\textbf{Claim:} $\kappa = \cos 2\theta_W = 7/13 \approx 0.538$.

\medskip
\textbf{Status: CONDITIONALLY VERIFIED.}

\medskip
\textbf{What works:}
\begin{itemize}
\item The spectral action on $\CP^2 \times S^3$ \emph{does} couple
$\psi$ asymmetrically to $U(1)_Y$ vs.\ $SU(2)_L$ (Check~1). \checkmark
\item The stiffness ratio $1/2$ determines $\sin^2\theta_W = 3/13$
via trace normalization (Check~2 prerequisite). \checkmark
\item $\kappa = 0$ is ruled out by internal consistency (Check~3). \checkmark
\item The formula $\cos 2\theta_W = 1 - 6/13 = 7/13$ is algebraically
exact and agrees with experiment to 0.17\% (Check~4). \checkmark
\item Assumptions 1--5 are internal to DFD (Check~5). \checkmark
\end{itemize}

\medskip
\textbf{What requires further work:}
\begin{itemize}
\item The identification of $\kappa$ with the photon mixing
asymmetry (Assumption~6) is a physically natural ansatz, not a
derived result.  It needs a rigorous derivation from the
gauge--$\psi$ Lagrangian on $\CP^2 \times S^3$.
\item The no-double-counting argument (stiffness sets $\theta_W$,
then mixing sets $\kappa$) is logically consistent but should be
formalized.
\item The tension with the Euler-Heisenberg route ($\kappa = 3/11$)
and hypercharge route ($\kappa = 1$) should be resolved by showing
these apply to different physical situations (photon-photon scattering
and bare charge counting, respectively, rather than $\psi$-photon
coupling).
\end{itemize}

\medskip
\textbf{No assumptions beyond DFD are required}, except for the
identification step (Assumption~6c) which, while natural, is not yet
derived from the DFD action.

\medskip
\textbf{The prediction $\kappa = 7/13$ is falsifiable} via the TE/TM
cavity frequency splitting:
\[
\frac{\delta\nu_{\mathrm{TE}} - \delta\nu_{\mathrm{TM}}}{\nu}
\approx 2\kappa\,\psi \approx 4.5 \times 10^{-6}.
\]
\end{tcolorbox}


%=============================================================================
\section{Derivation Chain Summary}
\label{sec:chain}
%=============================================================================

The full logical chain for $\kappa$:
\begin{equation}
\resizebox{0.9\columnwidth}{!}{$
\boxed{
\CP^2\!\times\! S^3
\xrightarrow{\text{Ricci}}
\kappa_r = n_r\kappa_0
\xrightarrow{\text{trace}}
\sin^2\!\theta_W = \frac{3}{13}
\xrightarrow{\text{mixing}}
\kappa = \frac{7}{13}
}
$}
\end{equation}

\begin{center}
\begin{tabular}{clc}
\toprule
\textbf{Step} & \textbf{Content} & \textbf{Status} \\
\midrule
1 & $\CP^2 \times S^3$ internal geometry & Theorem \\
2 & $\kappa_r = n_r\kappa_0$ (Frame Stiffness) & Theorem \\
3 & $\sin^2\theta_W = 3/13$ & Derived \\
4 & $\kappa = \cos 2\theta_W = 7/13$ & Ansatz \\
\bottomrule
\end{tabular}
\end{center}

Steps 1--3 are rigorous within DFD.  Step~4 is physically motivated
and self-consistent but remains an ansatz pending derivation from
the gauge--$\psi$ action.

\end{document}
